\subsection{Kết luận chương}

Chương 3 đã trình bày toàn cảnh kiến trúc bảo mật và các kỹ thuật tối ưu hiệu năng của hệ thống bỏ phiếu điện tử — một mô hình \textit{defense in depth} với bảy tầng phòng vệ độc lập, mỗi tầng vận hành một cơ chế kỹ thuật khác nhau và không phụ thuộc vào tầng khác để đảm bảo an toàn.

Về \textbf{trục bảo mật}, tầng mật mã EC-Schnorr ký mù tập thể đảm bảo tính bí mật và tính ẩn danh của phiếu bầu ngay cả trước máy chủ ký không trung thực; tầng JWT hai cặp khóa kết hợp xoay vòng refresh token chống tấn công đánh cắp phiên; phiên Redis với máy trạng thái và xóa nonce tức thời ngăn chặn cả bỏ phiếu trùng lẫn tấn công phục hồi khóa qua \textit{nonce reuse}; cây Merkle gắn với cam kết bất biến trên sổ cái Fabric mang lại khả năng kiểm chứng độc lập cho mọi cử tri và bên kiểm toán; mã hóa AES-256-GCM cho khóa bí mật của nút ký kết hợp \textit{authenticated encryption} bảo vệ tài sản nhạy cảm nhất ngay tại tầng lưu trữ tĩnh; cuối cùng tầng HTTP với Helmet, CORS, Rate Limiting, HTTPS và tầng TCP nội bộ với mTLS kết hợp \pkg{ValidationPipe} của \pkg{class-validator} cứng hóa toàn bộ bề mặt tấn công còn lại.

Về \textbf{trục hiệu năng}, hệ thống đạt thông lượng phù hợp cho quy mô bầu cử nhờ kết hợp nhiều kỹ thuật: cấu hình \texttt{BatchTimeout = 2s} ở tầng Orderer giúp Fabric gom được nhiều giao dịch vào cùng một khối, cân bằng giữa độ trễ và thông lượng; phân biệt nghiêm ngặt giữa \textit{invoke} và \textit{query} giúp tải đọc không cạnh tranh với tải ghi; cache hai tầng L1 LRU và L2 Redis giảm round-trip đến cơ sở dữ liệu; chiến lược chỉ mục MongoDB và pattern \textit{validate ngoài, commit trong giao dịch} đảm bảo throughput viết cao mà vẫn giữ tính nhất quán; truy vấn song song \pkg{Promise.all} ở mọi vị trí có thể giúp tổng độ trễ chỉ bằng độ trễ của thao tác chậm nhất.

Sáu tính chất của một hệ thống bầu cử an toàn — \textit{bí mật}, \textit{duy nhất}, \textit{xác thực}, \textit{không thể chối bỏ}, \textit{kiểm chứng được}, \textit{ẩn danh} — được đảm bảo đồng thời nhờ sự kết hợp của các tầng nêu trên, chứ không phải bởi một cơ chế duy nhất. Đây là điểm khác biệt căn bản giữa hệ thống đề xuất và các tiếp cận chỉ dựa vào blockchain công khai hay chỉ dựa vào chữ ký mù truyền thống.

Trên cơ sở thiết kế đã trình bày, Chương~4 tiếp theo sẽ trình bày quá trình hiện thực hóa, kết quả thử nghiệm trên môi trường mô phỏng và đánh giá hiệu năng thực tế của hệ thống dưới các kịch bản tải khác nhau.
